% Sorgente LaTeX canonico, importato e revisionato dal precedente sorgente Markdown.
\chapter{Confronto asintotico e simboli di Landau}
\label{chap:9}

\begin{obiettivocapitolo}{confrontare ordini di infinitesimo e infinito senza perdere termini decisivi}
\prioritaalta nei limiti e in Taylor. Ogni simbolo $o$, $O$ o $\sim$ deve avere un regime esplicito.
\end{obiettivocapitolo}
\begin{sceltametodo}{equivalenza oppure Taylor}
Usa $f\sim g$ quando l'espressione compare come fattore in un prodotto o in un quoziente e non vi è cancellazione. Usa Taylor quando compaiono somme/differenze di termini equivalenti, quando occorre il coefficiente del primo termine non nullo o quando serve una stima dell'errore.
\end{sceltametodo}
\begin{proceduraesame}{calcolo dell'ordine}
\begin{enumerate}
\item Porta tutte le funzioni allo stesso regime, spesso mediante $h=x-x_0$.
\item Sostituisci gli sviluppi fino all'ordine necessario.
\item Esegui prodotti, quozienti e differenze mantenendo i termini fino alla prima mancata cancellazione.
\item Raccogli la potenza dominante e concludi con equivalente, $o(\cdot)$ oppure $O(\cdot)$.
\end{enumerate}
\end{proceduraesame}
\begin{errorefrequente}{ordine insufficiente}
Se tutti i termini scritti si cancellano, non significa che il risultato sia zero: significa che lo sviluppo è stato arrestato troppo presto.
\end{errorefrequente}

Ogni scrittura asintotica deve essere accompagnata dal regime, per esempio \(x\to0\), \(x\to+\infty\) oppure \(n\to\infty\). Le equivalenze si possono sostituire in prodotti e quozienti, ma non in generale in somme o differenze nelle quali può verificarsi cancellazione.

Esempio:

\[
\sin x\sim x\quad(x\to0),
\]

ma sostituire direttamente \(\sin x\) con \(x\) in \(\sin x-x\) cancellerebbe il termine principale; serve uno sviluppo di ordine superiore:

\[
\sin x-x\sim-\frac{x^3}{6}.
\]

\section{Simboli di Landau}\label{simboli-di-landau}

Siano \(D\subseteq\mathbb R\), \(a\in\overline{\mathbb R}\) un punto di accumulazione nel regime considerato e \(f,g:D\to\mathbb R\). Gli insiemi \(U_\varepsilon\) e \(U\) indicano intorni del regime, puntati quando \(a\in\mathbb R\).

\[
\displaystyle
f=o(g)
\quad\Longleftrightarrow\quad
\forall\varepsilon>0\ \exists U_\varepsilon:
|f(x)|\le\varepsilon|g(x)|
\quad\forall x\in D\cap U_\varepsilon.
\]

Se \(g\) è definitivamente non nulla:

\[
\displaystyle
f=o(g)
\quad\Longleftrightarrow\quad
\frac{f}{g}\to0.
\]

\[
\displaystyle
f=O(g)
\quad\Longleftrightarrow\quad
\exists C>0\ \exists U:
|f(x)|\le C|g(x)|
\quad\forall x\in D\cap U.
\]

\[
\displaystyle
f=\Theta(g)
\quad\Longleftrightarrow\quad
\exists c,C>0\ \exists U:
c|g(x)|\le |f(x)|\le C|g(x)|
\quad\forall x\in D\cap U.
\]

Se \(g\) è definitivamente non nulla:

\[
\displaystyle
f\sim g
\quad\Longleftrightarrow\quad
\frac{f}{g}\to1
\quad\Longleftrightarrow\quad
f=g+o(g).
\]

Approfondimenti: \href{https://ingegnerismo.it/matematica/simboli-di-landau/}{simboli di Landau} ed \href{https://ingegnerismo.it/matematica/equivalenza-asintotica/}{equivalenza asintotica}.

Esercizi svolti: \href{https://ingegnerismo.it/matematica/simboli-di-landau-esercizi/}{simboli di Landau}.

\section{Algebra degli ordini}\label{algebra-degli-ordini}

Nel medesimo dominio e regime:

\[
\begin{aligned}
o(g)+o(g)&=o(g),&
O(g)+O(g)&=O(g),\\
f=o(g),\ g=O(h)&\Rightarrow f=o(h),&
f=O(g),\ g=o(h)&\Rightarrow f=o(h),\\
f=O(g),\ g=O(h)&\Rightarrow f=O(h),&
o(g)\,O(h)&=o(gh),\\
O(g)\,O(h)&=O(gh).&&
\end{aligned}
\]

\[
\displaystyle
g=\Theta(h)
\quad\Longrightarrow\quad
O(g)=O(h),
\qquad
o(g)=o(h).
\]

Se \(q\) è definitivamente non nulla:

\[
\displaystyle
r=o(g)
\quad\Longrightarrow\quad
\frac rq=o\!\left(\frac gq\right),
\]

e analogamente per \(O\).

\[
\displaystyle
\frac{a+o(1)}{b+o(1)}
=\frac ab+o(1),
\qquad b\ne0.
\]

Se \(f\sim g\) e \(h\sim k\), con denominatori definitivamente non nulli:

\[
\displaystyle
fh\sim gk,
\qquad
\frac fh\sim\frac gk.
\]

Approfondimento: \href{https://ingegnerismo.it/matematica/algebra-degli-o-piccoli/}{algebra degli o-piccoli}.

Esercizi svolti: \href{https://ingegnerismo.it/matematica/algebra-degli-o-piccoli-esercizi/}{algebra degli o-piccoli}.

\section{\texorpdfstring{Equivalenze fondamentali per \(x\to0\)}{Equivalenze fondamentali per x\textbackslash to0}}\label{equivalenze-fondamentali-per-xto0}

\[
\begin{gathered}
\sin x\sim x,
\qquad
\tan x\sim x,
\qquad
\arcsin x\sim x,
\qquad
\arctan x\sim x,\\[4pt]
1-\cos x\sim\frac{x^2}{2},\\[4pt]
e^x-1\sim x,
\qquad
\ln(1+x)\sim x,\\[4pt]
(1+x)^\alpha-1\sim\alpha x,
\qquad \alpha\ne0.
\end{gathered}
\]

\section{Gerarchie di crescita}\label{gerarchie-di-crescita}

Per \(x\to+\infty\), \(\alpha,\beta>0\) e \(a>1\):

\[
\displaystyle
(\ln x)^\beta=o(x^\alpha),
\qquad
x^\alpha=o(a^x).
\]

In notazione compatta:

\[
\displaystyle
(\ln x)^\beta\ll x^\alpha\ll a^x.
\]

Per \(n\to\infty\) e \(a>1\):

\[
\displaystyle
a^n\ll n!\ll n^n.
\]

Approfondimento: \href{https://ingegnerismo.it/matematica/gerarchia-degli-infiniti/}{gerarchia degli infiniti}.
